\documentclass[12pt]{article}
\usepackage[margin=1in]{geometry}
\usepackage[T1]{fontenc}
\usepackage{bold-extra}
\usepackage{xcolor}
\definecolor{garnet}{HTML}{73000A}

\setlength{\parindent}{0pt}
\setlength{\parskip}{0.3em}

\begin{document}

{\color{garnet}\large\textbf{\textsc{SC Space Grant Personal Essay}}}\\[-0.8em]
{\color{garnet}\rule{\textwidth}{0.5pt}}\\[0.2em]
I've always wanted to understand what's just out of view, like patterns beneath the surface or structure hidden in noise. That instinct sharpened into a question. How do you navigate a place where the sun never shines? Visible light won't work because there isn't any. GPS won't help since it doesn't exist on the Moon. The answer is multi-spectral perception, combining thermal cameras that see heat, radar that sees structure, and algorithms that fuse it all.

Throughout high school, I developed the discipline and teamwork that now drive my research. Four years of soccer and track demanded mental resilience, time management, and the ability to work within a team. Soccer taught me how individual contributions must align with collective goals, while track instilled personal accountability. Beyond athletics, HOSA showed me the importance of innovation in healthcare, while DECA provided insights into business strategy and entrepreneurship. These experiences broadened my perspective, showing me that meaningful solutions require not only technical knowledge but also an understanding of how systems operate within larger social and economic contexts.

Mechanical engineering at USC unified those instincts into a discipline. Freshman year, I sought research combining computation with physical systems. Dr.\ Wang's lab was developing multi-spectral perception for autonomous navigation using thermal cameras, visible sensors, and radar. I asked to join, not knowing Python would be the least of what I'd need to learn. Graduate students guided me through camera calibration, thermal-visible fusion, and deep learning. Now I'm co-authoring a conference paper on depth-aware registration, and the camera systems I helped build generate the data for LDAR.

Current thermal-visible registration methods face a tradeoff between assuming a flat scene and losing accuracy or adding LiDAR and losing simplicity. My research proposes a third path. Objects at different distances shift differently between cameras, a problem called parallax. The Learned Depth-Adaptive Registration framework trains neural networks to estimate depth from thermal imagery alone and adjust alignment accordingly, requiring no extra hardware. This makes robust multi-sensor fusion viable for platforms where every gram and watt matters.

NASA's Artemis missions will send rovers into the lunar South Pole, regions where sunlight never touches the surface and GPS doesn't exist. Navigation there demands sensors that work across extreme lighting including direct sun, deep shadow, and permanent darkness. Thermal imaging cuts through these conditions, but only if it aligns precisely with visible cameras. The techniques I'm developing apply directly to perception systems that will guide exploration of the Moon's permanently shadowed regions.

What draws me to this work is giving machines the ability to perceive what we cannot, such as heat signatures in shadow, terrain in permanent darkness, and obstacles obscured by dust. I aspire to become an engineer who bridges disciplines, drawing upon teamwork from athletics, strategic thinking from business, and a respect for science cultivated through health-focused studies. I plan to pursue graduate research in multi-modal sensing for space robotics, collaborating with professionals who share a commitment to advancing technology for the greater good.

\end{document}
